% БҮЛЭГ 8 — ДҮГНЭЛТ

\chapter{Дүгнэлт}
\label{Chapter8}

Энэхүү бүлэгт дипломын ажлын бүх бүлгийн оруулсан хувь нэмрийг нэгтгэн дүгнэж,
guideline-ийн заавал биелүүлбэл зохих шаардлагуудын биелэлтийг хүснэгтэнд
баримтжуулж, судалгааны хязгаарлалт, цаашдын чиглэл, нийгмийн нөлөөллийн
үнэлгээг танилцуулна.

%-------------------------------------------------------------------------------
\section{Бүлэг бүрийн оруулсан хувь нэмэр}

\textbf{Бүлэг 1 (Асуудлын томьёолол).} Монгол хэлний цахим орон зайд хортой
агуулга нэмэгдэж буй нөхцөл байдлыг тодорхойлж, судалгааны SMART зорилго,
зорилтуудыг томьёолсон. Судалгааны хамрах хүрээ (хоёр шатлалт ангилал, MN-BERT
суурьтай арга барил) болон үл хамаарах хязгаарлалтуудыг нарийвчлан зааглав.

\textbf{Бүлэг 2 (Холбогдох судалгааны тойм).} NLP, мэдрэмжийн шинжилгээ, хортой
агуулгын илрүүлэлт, Transformer/BERT архитектурын онолын үндсийг тоймлож,
Монгол хэл болон бусад агглютинатив хэлэнд хийгдсэн $\geq\!5$ судалгааг
харьцуулсан хүснэгт үүсгэв. Одоогийн орхигдсон чиглэл --- Монгол хэлний
Хортой сөрөг vs Бүтээлч шүүмжлэл ангиллын тусгай өгөгдлийн багц болон олон шатлалт системийн
дутагдал --- тодорхой болсон.

\textbf{Бүлэг 3 (Аргазүй ба архитектур).} Хоёр шатлалт ангиллын арга зүйг
сонгосон шалтгаан, MN-BERT нарийн тохируулгын стратеги, суурь загваруудын сонголтын
үндэслэлийг танилцуулав. Мөн AI хэрэгслийн ил тод ашиглалтыг guideline-ийн
зүйл 7-ын дагуу мэдүүлсэн.

\textbf{Бүлэг 4 (Өгөгдлийн боловсруулалт).} Найман нийтэд нээлттэй эх сурвалжаас
цуглуулсан Монгол кирилл шаардлага хангасан өгөгдлийг урьдчилсан боловсруулалтын
дамжлагаар дамжуулж, 4 ангиллын annotation стратегийг 3 шошгологч хэрэглэн
гүйцэтгэв. $7{,}043/1{,}428/1{,}529$ ($\approx$70.4/14.3/15.3\%) харьцаатай
train/val/test хуваалт хийж, inter-annotator agreement (Fleiss' $\kappa=0.72$)-ыг
тайлагнасан.

\textbf{Бүлэг 5 (Загвар хөгжүүлэлт ба дамжлага).} MN-BERT-ийн архитектурын
дэлгэрэнгүй, Stage~1 ба Stage~2-ын сургалтын дамжлага, гиперпараметрийн оновчлол,
ангиллын тэнцвэржүүлэлт (SMOTE суурь загварт; Focal Loss $\gamma{=}2.0$ + жин
$\alpha$ + WeightedRandomSampler эцсийн нэг шатлалт MN-BERT-д) болон Gradio-д
суурилсан инференсийн дамжлага-ыг баримтжуулав. Туршилтын давтагдах байдлын нөхцөлийг
random seed, хувилбарын хяналт, туршилтын лог бүртгэлээр хангасан.

\textbf{Бүлэг 6 (Үр дүн ба хэлэлцүүлэг).} Stage~1, Stage~2, end-to-end гурван
түвшинд тоон үр дүнг танилцуулж, MN-BERT нь TF-IDF суурьтай загвараас илүү
гүйцэтгэлтэй байх шалтгаануудыг тайлбарлав. Нэмэлт туршилт болгон
Transfer-init ablation аргыг хэрэгжүүлж, Stage~2 Macro F1 $0.9023 \to 0.9097$
болж сайжирсан. Focal Loss + Cosine Warmup-аар сургасан эцсийн нэг шатлалт загвар нь
нарийвчлал $\mathbf{83.26\%}$, Macro-F1 $\mathbf{0.8062}$-д хүрч, өмнөх
дипломын хамгийн сайн ажлын нарийвчлал ($82.70\%$) болон Macro-F1
($0.80$)-ыг тус тус давсан. Алдааны дүн шинжилгээ нь ёж, контекст хомсдол, богино текстэнд гардаг
гол алдааны хэв маягийг тодорхойлсон.

\textbf{Бүлэг 7 (Ёс зүйн асуудал).} Өгөгдлийн зөвшөөрөл, GDPR зарчмууд, annotation-ы
ёс зүй, bias ба тэгш байдал, нууцлал, хос хэрэглээний эрсдэл, AI хэрэгслийн
ашиглалт гэсэн зургаан хэмжээст ёс зүйн шинжилгээг хийж, guideline-ийн
Ш-5 шаардлагыг бүрэн хангасан.

\textbf{Бүлэг 8 (Дүгнэлт).} Энэ бүлэг --- бүх бүлгийн хувь нэмрийг нэгтгэж,
шаардлагын биелэлтийг хүснэгтээр харуулж, цаашдын судалгааны чиглэлийг
тодорхойлно.

%-------------------------------------------------------------------------------
\section{Guideline шаардлагын биелэлт}

Заавал биелүүлэх Ш-1 --- Ш-7 шаардлагыг үнэлгээний рубрикийн (guideline §10)
жишиг онооны хамт, мөн судалгааны чиглэлд тохирох сонгомол НШ шаардлагуудыг
дараах хүснэгтэд нэгтгэв.

\begin{table}[H]
\centering
\caption{Guideline шаардлагын биелэлтийн нэгдсэн хүснэгт (рубрикийн §10 дагуу)}
\label{tab:requirements}
{\small\setlength{\tabcolsep}{4pt}
\begin{tabular}{@{}p{0.9cm}p{3.1cm}p{5.5cm}p{2.7cm}@{}}
\toprule
\textbf{Код} & \textbf{Шаардлага} & \textbf{Биелэлтийн нотолгоо} & \textbf{Төлөв} \\
\midrule
Ш-1 & Загварын гүйцэтгэл & Бүлэг~\ref{Chapter6}: нэг шатлалт MN-BERT нарийвчлал $83.26\%$ ($80$--$89\%$ муж) & 3 оноо \\
\addlinespace[2pt]
Ш-2 & Өгөгдлийн чанарын сайжруулалт & Олон үе шаттай дахин шошгололт (Хүснэгт~\ref{tab:dq-progression}): нарийвчлал $0.7253{\to}0.8326$ ($\geq$10\%) & 4 оноо \\
\addlinespace[2pt]
Ш-3 & Архитектурын шинэлэг байдал & Нэг шатлалт vs хоёр шатлалт харьцуулсан судалгаа, MN-BERT нарийн тохируулга (Бүлэг~\ref{Chapter3},~\ref{Chapter5}) & 4 оноо \\
\addlinespace[2pt]
Ш-4 & Оновчлол ($\geq$20\% бууралт) & Загварын хөнгөвчлөл (distillation/ quantization) хийгдээгүй --- цаашдын ажил & Хэрэгжээгүй \\
\addlinespace[2pt]
Ш-5 & Ёс зүй, bias, privacy шинжилгээ & Бүлэг~\ref{Chapter7} бүхэлдээ & Бүрэн \\
\addlinespace[2pt]
Ш-6 & Дахин давтагдах туршилтын орчин & random seed, орчин, код, лог (Бүлэг~\ref{Chapter5}) & Давтагдахуйц \\
\addlinespace[2pt]
Ш-7 & Академик бичвэр ба хамгаалалт & Энэ бүхий л дипломын бичвэр, хамгаалалт & Сайн \\
\midrule
\multicolumn{4}{@{}p{12.2cm}@{}}{\textit{Сонгомол нэмэлт шаардлага (судалгааны чиглэлд тохирохыг сонгов):}}\\
\addlinespace[2pt]
НШ-1 & Суурь (baseline) загвар байгуулах & LR/NB/SVM суурь загвар, тоон утга тогтоосон (Бүлэг~\ref{Chapter6}) & $\checkmark$ \\
\addlinespace[2pt]
НШ-2 & $\geq\!3$ алгоритмын харьцуулалт & Хүснэгт~\ref{tab:comparison-summary} (4 загвар) & $\checkmark$ \\
\addlinespace[2pt]
НШ-11 & Абляцийн судалгаа & Transfer-init ablation (хэсэг~\ref{sec:transfer-init}), чанаржуулалтын явц (Хүснэгт~\ref{tab:dq-progression}) & $\checkmark$ \\
\addlinespace[2pt]
НШ-14 & Нийгмийн нөлөөллийн үнэлгээ & Нийгмийн нөлөөллийн үнэлгээ (Бүлэг~\ref{Chapter8}) & $\checkmark$ \\
\bottomrule
\end{tabular}}
\end{table}

\textit{Тэмдэглэл:} Рубрикийн §10-ийн дагуу Ш-1-ийн хувьд $80$--$89\%$ нь 3 оноо,
$\geq$90\% нь 4 оноо. Эцсийн нэг шатлалт MN-BERT загвар нь TF-IDF суурьтай
хамгийн сайн загвараас (нарийвчлал $66.91\%$) нарийвчлалаар $+16.35$ нэгжээр
давсан. Ш-4 (загварын хөнгөвчлөл) нь энэ ажлын хамрах хүрээнд багтаагүй бөгөөд
цаашдын судалгааны чиглэлд тодорхойлсон.

%-------------------------------------------------------------------------------
\section{Хязгаарлалт}

Судалгааны хязгаарлалтыг илэн далангүй баримтжуулах нь академик шударга байдлын
нэг илрэл юм. Дараах зургаан хязгаарлалтыг онцлов:

\begin{enumerate}
  \item \textbf{Өгөгдлийн цар хүрээ:} Одоогийн байдлаар 10,000 дээжтэй
    corpus нь Монгол хэлний цахим харилцааны бүх сэдвийн хамралтыг хараахан
    бүрэн илэрхийлээгүй. Эрүүл мэнд, санхүү, хууль шиг тусгай нэр томьёоны
    хэмжээнд өгөгдөл цөөн.
  \item \textbf{Платформын гажиг:} Найман эх сурвалж ($\approx$91\% мэдээний сайтын
    сэтгэгдэл) нь Монгол хэрэглэгчдийн зөвхөн хэсэг бүрэлдэхүүнийг төлөөлнө.
    TikTok, YouTube, Twitter/X гэх мэт бусад цахим платформын өгөгдөл
    судалгаанд ороогүй тул загварын ерөнхийлөх чадварт сөрөг нөлөө
    үзүүлж болзошгүй.
  \item \textbf{Ёж, контекст хамааралтай алдаа:} Бүлэг~\ref{Chapter6}-ын алдааны
    шинжилгээнд тодорхой харагдсанчлан, ёжтой илэрхийлэл, давхар утгатай
    үг, контекстээс хамаарсан аяс зэрэгт MN-BERT хэтэрхий өгүүлбэрийн
    гадаргуун шинжид найдаж ангилдаг.
  \item \textbf{Тооцоолох нөөцийн хязгаарлалт:} Сургалтыг локал AMD GPU орчинд
    явуулсан тул илүү том загвар (BERT-large, LLM нарийн тохируулга) туршиж
    амжаагүй. Мөн гиперпараметр хайлтыг бүрэн grid хайлтаар биш, тодорхой
    хүрээнд явуулсан.
  \item \textbf{Аннотацийн үлдэгдэл алдаа:} Датасетыг $v1$-ээс $v7$ хүртэл
    давтамжит чанаржуулалтаар --- загвараар тэмдэглэгдсэн өндөр итгэлтэй
    зөрүүг анхдагч шошгологчдын багаар дахин хянах замаар --- сайжруулсан
    боловч аннотацийн үлдэгдэл алдаа бүрэн арилаагүй байж болзошгүй. Том
    хэмжээний, хараат бус давтан шошгололт нь датасетын найдвартай байдлыг
    улам бататгана.
  \item \textbf{Үнэлгээний дисперс:} Үр дүнг нэг stratified hold-out хуваалт,
    нэг random seed ($42$) дээр тайлагнасан. Олон seed/k-fold cross-validation
    болон илүү өргөн статистик шинжилгээ хийх нь дисперсийн үнэлгээг
    сайжруулах боловч тооцоолох нөөцийн улмаас цаашдын ажил болгов.
\end{enumerate}

%-------------------------------------------------------------------------------
\section{Цаашдын судалгааны чиглэл}

Энэхүү ажлыг цаашид өргөжүүлэх чиглэлүүдийг доор жагсаав:

\begin{enumerate}
  \item \textbf{Өгөгдлийг өргөжүүлэх:} TikTok, YouTube коммент, Twitter/X-ийн
    Монгол хэлтэй контент, online форум зэрэг олон платформоос нэмэлт
    өгөгдөл цуглуулж, платформын олон талт байдлыг нэмэгдүүлэх. Ирээдүйд
    $\geq$50,000 шошготой дээж бүхий нэгдсэн Монгол хэлний toxicity corpus
    нийтийн эзэмшилд гаргахаар ажилламаар байна.
  \item \textbf{Илүү том загвар туршилт:} mBERT, XLM-R, LLaMA суурьтай
    Монгол хэлэнд дасан зохицуулсан pre-train загварууд, эсвэл GPT
    төрлийн causal LM-ийг few-shot/instruction нарийн тохируулга аргаар
    ашиглаж гүйцэтгэлийн дээд хязгаарыг хэмжих.
  \item \textbf{Ёж, контекст ойлгогч модуль:} Ёжлолт илрүүлэлтийг (sarcasm detection) нэмэлт
    тусдаа даалгавар болгон multi-task learning хэлбэрээр нэгтгэх. Мөн
    коммент нь хариу өгч буй эх нийтлэл эсвэл дээд коммент руу хамтад нь
    ачаалах context-aware архитектур судлах.
  \item \textbf{Тайлбарлах чадвар (Explainability):} SHAP, LIME, attention
    визуализаци ашиглан хэрэглэгчид яагаад тухайн ангилал хүлээн авсан
    болохыг тайлбарлах модуль нэмэх. Энэ нь ёс зүйн ``human-in-the-loop''
    загварт зайлшгүй чухал.
  \item \textbf{Загварын хөнгөвчлөл (Distillation \& Quantization):}
    Инференсийн хурдыг нэмэгдүүлэх, мобайл/edge үйлчилгээнд ашиглах
    зорилгоор knowledge distillation, INT8 quantization туршилтыг
    хийх. Энэ нь бодит үйлдвэрлэлд хэрэглэх боломжийг нэмэгдүүлнэ.
  \item \textbf{End-to-End нарийн тохируулга (Joint fine-tuning):}
    Хэсэг~\ref{sec:transfer-init}-ийн Transfer-init туршилтын үр дүнд
    үндэслэн Stage~1 ба Stage~2-ыг тусад нь биш, хамтад нь нэгдсэн
    pipeline хэлбэрээр end-to-end нарийн тохируулах стратеги судлах.
    Cascade алдааны дамжуулалтыг бууруулах боломжтой бөгөөд нэгдсэн
    системийн Macro F1-ийг нэмэгдүүлэх суурь болж чадна.
\end{enumerate}

%-------------------------------------------------------------------------------
\section{Нийгмийн нөлөөллийн үнэлгээ}

Монгол хэлний цахим орон зайд хортой агуулга нэмэгдэх нь хувь хүн, бүлэг,
институци тус бүрд сэтгэл зүй, нийгэм-улс төрийн сөрөг үр дагавар үзүүлдэг
онцгой асуудал юм. Энэхүү системийг хариуцлагатай, ил тод хэрэглэх нөхцөлд
нийгэмд дараах эерэг нөлөө үзүүлэх боломжтой: (1)~сэтгэл хоолой, онлайн
цагдан мөрдлөг гэх мэт хор уршгийн түвшинг урьдчилж илрүүлэх, (2)~цахим
платформууд нь өөрсдөө хяналт тавихад Монгол хэлний туслах хэрэгсэлтэй
болох, (3)~бүтээлч шүүмжлэлийг тодорхой ялган хамгаалснаар ардчилсан
хэлэлцүүлэг, иргэдийн оролцоог сурталчлах. Гэсэн хэдий ч систем нь Ш-5-д
дурдсанчлан зөвхөн ``зөвлөгч'' үүрэгтэй байх ёстой ба шийдвэрийг хүн
гаргах эцсийн хариуцлагатай; учир нь алгоритмын зохицуулалт нь оршин тогтнох,
үзэл бодлоо илэрхийлэх эрхтэй шууд холбогддог тул ардчилсан нийгмийн дархлааг
хамгаалах хариуцлага нь хүний гарт бэхлэгдсэн хэвээр байх ёстой.

%-------------------------------------------------------------------------------
\section{Бүлгийн дүгнэлт}

Энэ бүлэгт дипломын ажлын найман бүлгийн хувь нэмрийг нэгтгэж, guideline-ийн
Ш-1...Ш-7 заавал шаардлага болон сонгомол НШ шаардлагуудын биелэлтийг
үнэлгээний рубрикийн дагуу хүснэгтэнд баримтжуулав. Судалгааны зургаан гол хязгаарлалтыг тодорхойлов: өгөгдлийн цар хүрээ, платформын
гажиг, ёжийн алдаа, тооцоолох нөөц, аннотацийн үлдэгдэл алдаа, үнэлгээний дисперс. Мөн таван цаашдын судалгааны чиглэлийг тогтоов:
өгөгдөл өргөтгөх, том загвар туршилт, ёж илрүүлэх модуль, тайлбарлах чадвар,
загварын хөнгөвчлөл. Нийгмийн
нөлөөллийн үнэлгээнд энэхүү системийг зөвхөн хүний хяналтын дор зөвлөгчийн
үүрэгт ашиглах шаардлагатайг онцолсон. Эцсийн нэгдсэн үр дүн: нэг шатлалт MN-BERT
загварын нарийвчлал $\mathbf{83.26\%}$, Macro-F1 $\mathbf{0.8062}$ --- өмнөх
дипломын хамгийн сайн ажлын нарийвчлал ($82.70\%$) ба Macro-F1
($0.80$)-ыг тус тус давсан. Ийнхүү дипломын ажил нь Монгол хэлний
NLP хэрэглээнд чухал хувь нэмэр оруулж байгаа бөгөөд уг салбарын цаашдын
хөгжилд тодорхой чиглэл үзүүлж байна.
